\scanpage{177}
% L01–07
% Three diagonal strokes cross the opening draft.
\deleted{Donc \unsure{pr|il} \deleted{\unsure{choisir|s'agit d'}} exprimer l'axiome (GH1) du (3) (p.~81), mais en y prenant $f\in G$, de \unsure{façon précise, aussi} \deleted{Il est} clair qu'on doit pouvoir prendre $=$ $f$ \unsure{$\in$} provenant d'un élément de $B'$ (où $B'$ opère sur $S\simeq G/\widetilde U$ via son opération sur $(G\supset\widetilde U)$). Donc cette condition devient (dans la terminologie de (26)}
% slanted
\margin{\deleted{\unsure{qualifié|grâce} ``cartogr'' \unsure{par rapport à} $G_0$}}

% L08–13
\deleted{$\forall b''\in\widetilde B''$, définissant donc une permutation $g\widetilde U\longmapsto g\unsure{b''|b'}\widetilde U$ de $S\simeq G_0/\widetilde U$, qui commute aux op. de $G_0$ (et $=$ de $G$), il existe un $b'\in B'$ (dépendant de $b''$) tel que pour tout couple $(s,t)\in S$, on ait}

\[
\deleted{\{s,t\}\in A_0\Longleftrightarrow
\{b'^{-1}s,b'^{-1}\cdot b''t\}\in A_0} \tag*{\deleted{(58)}}
\]
\deleted{\illegible{1} on a \unsure{$s\in g$|$s=g$}}

% L14–15
\[
\begin{aligned}
&\deleted{\text{Or }\{s,t\}\in A_0\Longleftrightarrow\exists u\in G_0,\quad s=u(s_0),\quad t=u(t_0)}\\
&\deleted{\{b'^{-1}s,b'^{-1}\cdot b''t\}\in A_0\Longleftrightarrow\exists v\in G_0,\quad b'^{-1}s=vs_0,\quad b'^{-1}b''t=vt_0}
\end{aligned}
\]

% L16–18
\deleted{donc il faut écrire}
\[
\left\{\begin{aligned}
&\deleted{\forall\unsure{\lambda|b''}\in\unsure{Z|\widetilde B''}\subset G_0,\quad\exists f\in G\ \text{\unsure{(en réserve,} } f\in B')},\\
&\deleted{\text{tels que, }\forall s,t\in S,\ \text{on ait}}\\
&\deleted{\{s,t\}\in A_0\Longleftrightarrow\{f^{-1}s,f^{-1}\lambda t\}\in A_0}
\end{aligned}\right. \tag*{\deleted{(58)}}
\]

% L19–23
\deleted{On voit d'abord que pour $\lambda,f$ fixés, \inserted{pour vérifier la condition} du fait que $\lambda$ commute à $G$, il suffit de la vérifier pour $s=s_0$ \unsure{(quitte} et $t$ arbitraire, car si $s=us_0$ ($u\in G_0$), \unsure{on suppose $t=u(t')$} \unsure{il suffit où on|on a}}
\[
\begin{aligned}
&\deleted{\{s,t\}\in A_0\Longleftrightarrow\{\underbrace{u^{-1}s,u^{-1}t}_{s_0,\ u^{-1}t}\}\in A_0}\\
&\deleted{\{f^{-1}s,f^{-1}\lambda t\}\in A_0\Longleftrightarrow
\{u^{-1}f^{-1}\uline{u}u^{-1}s,\ u^{-1}f^{-1}uu^{-1}\lambda t\}\in A_0}\\
&\hspace{9em}\deleted{\{u^{-1}(f^{-1})s_0,\ u^{-1}(f^{-1})\lambda\uline{u^{-1}t}\in A_0}
\end{aligned}
\]

% L24–27
Il faut exprimer que si $\lambda\in Z\subset G$ et $f\in G$ sont tels que l'image de $f$ dans $G/G^0=Z'$ soit \deleted{\unsure{est}} égale à $\alpha(\lambda)$, \unsure{où|via} $\alpha:Z\simeq Z'$ (57 d), alors on a la relation

% archive sheet number: 176
